\documentclass[11pt]{article}

\usepackage[letterpaper,margin=0.78in]{geometry}
\usepackage[T1]{fontenc}
\usepackage[utf8]{inputenc}
\usepackage{lmodern}
\usepackage{microtype}
\usepackage{booktabs}
\usepackage{array}
\usepackage{amsmath}
\usepackage{graphicx}
\usepackage{enumitem}
\usepackage{seqsplit}
\usepackage[table]{xcolor}
\usepackage{hyperref}

\definecolor{passgreen}{HTML}{166534}
\definecolor{failred}{HTML}{B91C1C}
\definecolor{noteblue}{HTML}{1D4ED8}
\newcommand{\pass}{\textcolor{passgreen}{\textbf{PASS}}}
\newcommand{\fail}{\textcolor{failred}{\textbf{FAIL}}}
\newcommand{\code}[1]{\texttt{#1}}
\newcommand{\hash}[1]{{\ttfamily\seqsplit{#1}}}
\newcommand{\yes}{\textcolor{passgreen}{Yes}}
\newcommand{\no}{\textcolor{failred}{No}}

\hypersetup{
  colorlinks=true,
  linkcolor=noteblue,
  urlcolor=noteblue,
  pdftitle={E0.1 200-Tick Variable-Agent Compatibility Supplement},
  pdfauthor={AlemDICE Research Team}
}
\setlist{nosep,leftmargin=1.4em}
\setlength{\parindent}{0pt}
\setlength{\parskip}{0.55em}
\setlength{\emergencystretch}{3em}

\title{\textbf{E0.1: 200-Tick Variable-Agent Compatibility}\\
\large Engineering Supplement for the Agent-Scaling Pipeline}
\author{AlemDICE Research Team}
\date{Run completed 23 July 2026}

\begin{document}
\maketitle

\begin{abstract}
E0.1 stress-tested the unchanged Alem evaluator at
$N\in\{1,2,3,4,6\}$ with two seeds and a requested 200-tick horizon.
Provider-free scripted agents selected \code{Noop} and emitted deterministic
broadcasts, isolating population-dependent infrastructure from language-model
behavior.  All ten episodes produced complete artifacts, and every structural
gate passed: indexing, trajectory axes, targeted \code{Give} mappings,
specialization cycling, broadcast accounting, action parsing, and deterministic
saved-state reconstruction.  The strict aggregate gate nevertheless
\fail{} because three episodes ended naturally when all agents died before
tick 200: $N=1$/seed 13000 at tick 114, $N=4$/seed 13001 at tick 182, and
$N=6$/seed 13001 at tick 192.  Thus E0.1 clears the engineering compatibility
blocker for all tested populations, but it does not show that every episode can
reach a fixed horizon and provides no evidence about behavioral performance or
coordination quality.
\end{abstract}

\section{Question and decision rule}

The purpose of E0.1 was to determine whether the canonical evaluation and
artifact pipeline remains mechanically valid as the number of physical agents
changes.  It was not a performance experiment.  The requested matrix comprised
five population sizes, seeds 13000 and 13001, and at most 200 environment ticks
per episode, for ten episodes total.

Two decisions are reported separately:

\begin{description}[style=nextline,leftmargin=1.5em]
  \item[Structural gate.]
  An episode may terminate naturally, but must produce a complete canonical
  artifact and pass every shape, action, role, communication, replay, and
  provider-accounting check.
  \item[Strict horizon gate.]
  In addition to the structural requirements, every episode must execute all
  200 requested ticks.  In the runner's published summary,
  \code{passed} is the conjunction that includes this requirement.
\end{description}

This distinction is essential: an \code{environment\_terminated} episode is a
valid terminal trajectory, but it fails a protocol that literally requires 200
executed ticks.

\section{Methods}

\subsection{Environment and execution}

The run used \code{Alem-Coop-Symbolic}, Easy coordination difficulty,
\code{god\_mode=false}, soft specialization, unshared rewards, and the
\code{baseline} communication topology.  Each population used exactly one
indexed fake client configuration per physical agent, one episode worker, and
no debriefs or saved images.  Populations ran sequentially.  The evaluator's
normal per-tick worker pool still scheduled one actor call per physical agent.

The specialization cycle was Warrior, Forager, Miner, repeated by agent index;
its numeric representation was $(2,1,3)$.  Consequently, the expected cycles
were $(2)$, $(2,1)$, $(2,1,3)$, $(2,1,3,2)$, and
$(2,1,3,2,1,3)$.

\subsection{Scripted policy and communication}

At scheduled actor turn $t$, agent $i$ returned the valid action
\code{Noop} and the deterministic payload
\[
\text{\code{E0|AGENT=}}\;i\;\text{\code{|TICK=}}\;t\;
\text{\code{|STATE=ACTIVE}}.
\]
The response declared model ID \code{e0-scripted} and zero input, output, and
reasoning tokens.  No provider request was made.  Under baseline broadcast, the
expected accounting for an episode of $T$ executed ticks was
\[
M_{\mathrm{emit}}=NT,\qquad
M_{\mathrm{deliver}}=N(N-1)T.
\]
For $N=1$, messages were recorded as emitted but had no recipient and therefore
zero deliveries and zero delivered bytes.

\subsection{Structural validation}

For every population and seed, the harness checked:

\begin{itemize}
  \item complete episode JSON, trajectory, debug, and saved-state artifacts;
  \item observation, action, reward, text, and trajectory axes consistent with
        the configured $N$ and executed horizon;
  \item all ordered \code{Give to Agent $j$} mappings for $i\ne j$, including
        canonical round trips and action-slot indices;
  \item the expected specialization cycle;
  \item exact emitted-message and fan-out counts;
  \item action parse rate equal to 1.0;
  \item zero model calls, provider requests, and transport errors; and
  \item equality between a stable projection of the first saved state and a
        fresh reset at the same seed.  The projection hashes the map, item map,
        player positions, player specializations, and coordination map.
\end{itemize}

\subsection{Timing definitions}

All timing values use a monotonic clock.
\code{episode\_wall\_seconds} begins immediately before the main evaluator loop
and ends after the final tick, before most post-episode artifact processing.
\code{mean\_tick\_wall\_seconds} measures from the start of a tick through
incoming-message injection, state snapshots, concurrent actor calls, action
parsing, broadcast routing, and \code{env.step}; population means are weighted
by executed ticks.  \code{max\_tick\_wall\_seconds} is the largest such tick in
either seed.  \code{worker\_round\_wall\_seconds} covers only submission and
settlement of the concurrent scripted actor futures.  Total run elapsed time
also includes environment construction and compilation, reset/replay
validation, serialization, and between-population overhead.

Because every episode contains one much larger maximum, the results also report
a descriptive sensitivity statistic: accumulated tick time after removing
exactly one maximum from each episode, divided by the remaining ticks.  This is
labeled ``mean excluding per-episode max.''  It is derived from the recorded
means, maxima, and tick counts; it is not a separately timed warm or
steady-state benchmark.  ``Worker/tick'' divides cumulative worker-round time by
executed ticks.

\section{Results}

\subsection{Gate outcomes}

\begin{table}[ht]
\centering
\small
\begin{tabular}{@{}p{0.68\linewidth}c@{}}
\toprule
Published gate & Outcome\\
\midrule
All populations produced complete artifacts & \yes\\
All structural checks passed & \yes\\
All requested 200-tick horizons were reached & \no\\
Action parse rate was 1.0 in every episode & \yes\\
Zero model calls / provider requests / transport errors & \yes\\
Trajectory agent axes matched configured populations & \yes\\
All targeted \code{Give} mappings round-tripped & \yes\\
Specialization cycles matched & \yes\\
Broadcast accounting was exact & \yes\\
Saved-state replay projections matched & \yes\\
\midrule
\textbf{Structural decision} & \pass\\
\textbf{Strict decision, including full horizon} & \fail\\
\bottomrule
\end{tabular}
\caption{The strict result is false solely because the full-horizon gate is
false.  No structural or transport gate failed.}
\label{tab:gates}
\end{table}

\subsection{Population summary}

\begin{table}[ht]
\centering
\small
\resizebox{\textwidth}{!}{%
\begin{tabular}{@{}rrrrrrrrrcc@{}}
\toprule
$N$ & Eps. & Ticks & Agent-ticks & Obs. dim. & Actions & Give maps &
Deaths & Full horizon & Structural & Strict\\
\midrule
1 & 2 & 314 & 314  & 9,476  & 55 & 0  & 1 & No  & \pass & \fail\\
2 & 2 & 400 & 800  & 9,597  & 55 & 2  & 1 & Yes & \pass & \pass\\
3 & 2 & 400 & 1,200& 9,718  & 56 & 6  & 0 & Yes & \pass & \pass\\
4 & 2 & 382 & 1,528& 9,839  & 57 & 12 & 7 & No  & \pass & \fail\\
6 & 2 & 392 & 2,352& 10,081 & 59 & 30 & 8 & No  & \pass & \fail\\
\midrule
All & 10 & 1,888 & 6,194 & --- & --- & 50 & 17 & No & \pass & \fail\\
\bottomrule
\end{tabular}}
\caption{Population-level compatibility results.  ``Give maps'' counts ordered
giver--recipient mappings validated once per population, hence $N(N-1)$.}
\label{tab:population}
\end{table}

Observation length increased by 121 values per added agent:
\[
d_{\mathrm{obs}}(N)=9476+121(N-1).
\]
The action surface was 55 actions for $N=1$ and $N=2$, then increased by one
targeted handoff slot per additional peer, reaching 59 at $N=6$.  These are
interface measurements, not estimates of behavioral benefit.

\subsection{Seed-level termination and mortality}

\begin{table}[ht]
\centering
\scriptsize
\resizebox{\textwidth}{!}{%
\begin{tabular}{@{}rrrrlrrrrrr@{}}
\toprule
$N$ & Seed & Ticks & Deaths & Terminal code & Emit & Deliver & Bytes &
Episode wall (s) & Mean tick (s) & Max tick (s)\\
\midrule
1 & 13000 & 114 & 1 & \code{environment\_terminated} & 114 & 0 & 0
  & 20.7634 & 0.1817 & 11.9478\\
1 & 13001 & 200 & 0 & \code{environment\_truncated}  & 200 & 0 & 0
  & 26.8550 & 0.1340 & 11.7007\\
2 & 13000 & 200 & 1 & \code{environment\_truncated}  & 400 & 400 & 12,580
  & 34.4896 & 0.1719 & 10.1435\\
2 & 13001 & 200 & 0 & \code{environment\_truncated}  & 400 & 400 & 12,580
  & 35.9126 & 0.1791 & 10.5244\\
3 & 13000 & 200 & 0 & \code{environment\_truncated}  & 600 & 1,200 & 37,740
  & 47.2818 & 0.2356 & 11.4559\\
3 & 13001 & 200 & 0 & \code{environment\_truncated}  & 600 & 1,200 & 37,740
  & 47.3205 & 0.2361 & 11.4955\\
4 & 13000 & 200 & 3 & \code{environment\_truncated}  & 800 & 2,400 & 75,480
  & 58.4060 & 0.2912 & 11.2831\\
4 & 13001 & 182 & 4 & \code{environment\_terminated} & 728 & 2,184 & 68,568
  & 55.1834 & 0.3025 & 12.0740\\
6 & 13000 & 200 & 2 & \code{environment\_truncated}  & 1,200 & 6,000 & 188,700
  & 86.5491 & 0.4318 & 11.9992\\
6 & 13001 & 192 & 6 & \code{environment\_terminated} & 1,152 & 5,760 & 181,020
  & 85.4465 & 0.4439 & 11.7557\\
\bottomrule
\end{tabular}}
\caption{Exact per-seed results.  \code{environment\_truncated} denotes reaching
the configured 200-tick cap; \code{environment\_terminated} denotes natural
termination after all agents died.}
\label{tab:seeds}
\end{table}

All 17 deaths were attributed to mob combat.  There were zero deaths from
starvation, dehydration, exhaustion, or friendly fire.  Partial mortality did
not terminate an episode: $N=2$/seed 13000 reached the cap with one death,
$N=4$/seed 13000 with three deaths, and $N=6$/seed 13000 with two deaths.
Natural termination occurred only in the three rows where the death count
equaled $N$.

\subsection{Communication accounting}

\begin{table}[ht]
\centering
\small
\begin{tabular}{@{}rrrrrr@{}}
\toprule
$N$ & Executed ticks & Emitted & Delivered & Delivery bytes &
Deliveries / emission\\
\midrule
1 & 314 & 314   & 0      & 0       & 0\\
2 & 400 & 800   & 800    & 25,160  & 1\\
3 & 400 & 1,200 & 2,400  & 75,480  & 2\\
4 & 382 & 1,528 & 4,584  & 144,048 & 3\\
6 & 392 & 2,352 & 11,760 & 369,720 & 5\\
\midrule
All & 1,888 & 6,194 & 19,544 & 614,408 & ---\\
\bottomrule
\end{tabular}
\caption{Broadcast accounting across both seeds.  Every row exactly matches
$NT$ emissions and $N(N-1)T$ deliveries.}
\label{tab:communication}
\end{table}

The result makes the baseline's quadratic delivery surface explicit.  At
$N=6$, 2,352 scripted emissions produced 11,760 delivered copies.  Because the
probe emits on every scheduled configured-agent turn, these counts remain
algebraic even after individual agents die; they validate accounting, not useful
information flow among living agents.

\subsection{Timing}

\begin{table}[ht]
\centering
\scriptsize
\resizebox{\textwidth}{!}{%
\begin{tabular}{@{}rrrrrrr@{}}
\toprule
$N$ & Episode wall total (s) & Mean episode (s) & Weighted mean tick (s) &
Mean excl.\ per-episode max (s) & Max tick (s) & Worker/tick (ms)\\
\midrule
1 & 47.6183  & 23.8092 & 0.151301 & 0.076475 & 11.9478 & 0.155216\\
2 & 70.4022  & 35.2011 & 0.175483 & 0.124436 & 10.5244 & 0.176712\\
3 & 94.6023  & 47.3012 & 0.235863 & 0.179381 & 11.4955 & 0.165132\\
4 & 113.5894 & 56.7947 & 0.296600 & 0.236695 & 12.0740 & 0.176943\\
6 & 171.9955 & 85.9978 & 0.437702 & 0.379037 & 11.9992 & 0.186564\\
\bottomrule
\end{tabular}}
\caption{Population timing from the raw episode artifacts.  Means are weighted
by executed ticks; the exclusion statistic removes one maximum from each seed.}
\label{tab:timing}
\end{table}

The ten episode timers summed to 498.207821 seconds.  End-to-end run elapsed
time was 659.272638 seconds, leaving 161.064817 seconds in construction,
compilation, reset/replay validation, serialization, and orchestration outside
the episode timers.  The tick means increase with $N$, while the summed
scripted worker-future time was only 0.326201 seconds across the entire matrix.
The roughly 10--12 second maximum tick at every population is consistent with
one-time compiled-environment overhead inside the measured loop; E0.1 was not
designed as a warmed throughput benchmark.  The artifacts store only the mean
and maximum, not the maximum's tick index or cause, so this interpretation is
not a proven attribution.  The mean-excluding-maximum column ranges from
0.076475 seconds at $N=1$ to 0.379037 seconds at $N=6$, but no exact
warm/steady-state timing is reported.

\section{Interpretation and promotion decision}

\textbf{Engineering conclusion.}  The structural gate is \pass{} at all five
population sizes.  E0.1 found no indexing, action-space, trajectory,
serialization, broadcast-accounting, replay, or provider-isolation defect.
Therefore $N=1,2,3,4$ may proceed to the supported-range Source screen, and
$N=6$ clears the engineering prerequisite for a separately labeled extension.

\textbf{Protocol conclusion.}  The strict 200-tick gate is \fail{}.  A future
report must not describe this matrix as ``ten 200-tick episodes'': seven
episodes reached the cap, while three ended naturally at 114, 182, and 192
ticks.  For behavioral experiments, natural all-agent death should remain a
valid outcome and be analyzed as mortality and time-to-termination rather than
misclassified as an API or infrastructure failure.  Comparisons must normalize
cost and opportunity by actual executed agent-ticks as well as report the
requested horizon.

\textbf{Scientific claim boundary.}  The scripted \code{Noop} policy cannot
show whether adding agents improves return, achievements, cooperation,
communication efficiency, or robustness.  Mortality patterns in two seeds are
diagnostic only and must not be interpreted as a population-size effect.

\section{Limitations}

\begin{itemize}
  \item Only two seeds were used, with a deliberately inactive policy.
  \item No hosted model, reasoning, prompt-context growth, or generated-message
        parsing was exercised.
  \item Broadcasts were deterministic synthetic payloads; semantic value,
        redundancy, and teammate response were not tested.
  \item The fixed 48-by-48 world, spawn geometry, mob pressure, role balance,
        and coordination opportunities are not controlled independently of
        population.
  \item $N=1$ uses the solo wrapper, and $N=6$ lies beyond the paper's
        documented one-to-four range; both require explicit labels in later
        figures.
  \item Timing includes cold compiled-environment effects and is unsuitable for
        estimating hosted-model throughput or cost.
  \item The deterministic replay check covers the stated stable state
        projection, not bitwise identity of every field or a full re-executed
        trajectory.
\end{itemize}

\section{Provenance and audit trail}

{\raggedright
The run began at \code{2026-07-23T06:57:25.963656+00:00}, ended at
\code{2026-07-23T07:08:25.236294+00:00}, and ran for 659.272638 seconds.
It used Python 3.12.13 on
\path{Linux-6.17.0-40-generic-x86_64-with-glibc2.42}.  The source branch was
\path{research/agent-scaling-recruitment-pipeline} at commit
\hash{5c4c64338b9aae9132f9cc32b6eec3d746964529}.  The source-status record
contained one pre-existing untracked replay,
\path{Results/replays/nano_high_source_full_world.mp4}; it was outside the run
root and was not an input to this report.

The canonical run root is
\path{outputs/alem_eval/e0_1_agent_compatibility_200_v2}.  It contains 82 files
occupying 78,221,137 bytes.  The source-of-truth summary
\path{e0_results.json} has SHA-256
\hash{6d28e885e241b4f58553291d6fa205fb982878f2575f6bdd3a24d711fda9b529}.
The committed runner at the recorded source revision has SHA-256
\hash{b2aeaca114686b547fe3b205351d3e483c4576f391194bc744dea06ed1d806a9}.
\par}

Resolved configuration SHA-256 values were:
\begin{center}
\small
\begin{tabular}{@{}rl@{}}
\toprule
$N$ & \code{resolved\_config.yaml} SHA-256\\
\midrule
1 & \hash{2ed69835be2ef2a8e4779221e0c54f06d65cf10b48f18ca3406e87ed42762160}\\
2 & \hash{7895941037e47287b0fe671952dcc57fca81d66fb4d7a61edea3592338bda619}\\
3 & \hash{815c7d344cb690af1a5ba396838000c6c64a6ef72c032be88ed1e4ec3d393ca0}\\
4 & \hash{d0fbf4db8f0252c5ffc28f2fedff20345c113875bcd169f913e5b4638e17b41a}\\
6 & \hash{93583d0d834bf7fdc1226505b99c7e0e837dd028944ffe1202efc49300f669c6}\\
\bottomrule
\end{tabular}
\end{center}

A corresponding reproducible invocation is:
\begin{quote}
\small\ttfamily
uv run --extra baselines-llm --python 3.12\\
\hspace*{1em}python scripts/run\_e0\_agent\_compatibility.py\\
\hspace*{1em}--output outputs/alem\_eval/e0\_1\_agent\_compatibility\_200\_v2\\
\hspace*{1em}--counts 1,2,3,4,6 --seeds 13000,13001 --steps 200
\end{quote}

\section{Conclusion}

E0.1 validates variable-agent infrastructure through six agents under a
provider-free 200-tick stress protocol.  Its correct result is not a single
undifferentiated pass or fail: the structural compatibility decision is
\pass{}, while the literal full-horizon decision is \fail{} because all-agent
death naturally ended three trajectories.  This evidence supports proceeding
to carefully normalized behavioral screens, with deaths and actual agent-ticks
treated as first-class outcomes.

\end{document}
